% =====================================================================
% anexo.tex -- Anexo: metodologia de elaboracao deste esboco.
% Mapeia, secao a secao, quais partes do texto derivam da NUREG-1537,
% Parte 1 ("Format and Content") e Parte 2 ("Standard Review Plan and
% Acceptance Criteria"), e quais elementos sao sintese propria, sem
% base direta na norma.
% =====================================================================
\nusecbreak{anexo1}{Anexo}
\section*{Anexo -- Metodologia de Elabora\c{c}\~ao deste Esbo\c{c}o}
\addcontentsline{toc}{section}{Anexo -- Metodologia de Elabora\c{c}\~ao deste Esbo\c{c}o}
\label{sec:anexo-metodologia}

Este anexo documenta, para cada se\c{c}\~ao do Cap\'itulo~1, a
correspond\^encia entre o texto deste esbo\c{c}o e as duas partes da
NUREG-1537 \citep{nureg1537p1,nureg1537p2}, distinguindo o que \'e
adapta\c{c}\~ao direta da norma do que \'e s\'intese pr\'opria,
incorporada apenas como ferramenta de revis\~ao interna do requerente.

\textbf{Regra geral de composi\c{c}\~ao de cada se\c{c}\~ao (1.1 a 1.8):}

\begin{itemize}
\item O \textbf{texto corrido} segue a NUREG-1537, Parte~1
  (\emph{``Format and Content''}) \citep{nureg1537p1}, se\c{c}\~ao
  correspondente -- \'e uma adapta\c{c}\~ao/tradu\c{c}\~ao do conte\'udo
  m\'inimo ali exigido, com os colchetes \texttt{[...]} substituindo os
  dados espec\'ificos do requerente.
\item A \textbf{tabela de requisitos de aceita\c{c}\~ao} reproduz, em
  ordem e conte\'udo, os itens da subse\c{c}\~ao \emph{``Acceptance
  Criteria''} da NUREG-1537, Parte~2 (\emph{``Standard Review Plan and
  Acceptance Criteria''}) \citep{nureg1537p2}, se\c{c}\~ao
  correspondente -- \'e a lista que o revisor do \'org\~ao regulador
  efetivamente usa para avaliar a se\c{c}\~ao.
\item A \textbf{tabela-resumo dos pontos mais importantes} e a
  \textbf{tabela de ap\^endices/refer\^encias cruzadas a verificar}
  \textbf{n\~ao existem na NUREG-1537}. S\~ao s\'intese pr\'opria,
  acrescentada como ferramenta de revis\~ao interna -- \'util para o
  requerente conferir consist\^encia entre este cap\'itulo-resumo e os
  cap\'itulos detalhados do RPAS, mas n\~ao devem ser citadas como se
  fossem conte\'udo da norma.
\end{itemize}

\begin{tabelametodologia}
1.1 &
  Parte 1, \S1.1 \citep{nureg1537p1} &
  Parte 2, \S1.1, \emph{Acceptance Criteria} (6 itens) \citep{nureg1537p2} &
  Tabela-resumo; tabela de ap\^endices/refer\^encias cruzadas (Cap.~2--17
  do RPAS) \\
1.2 &
  Parte 1, \S1.2 \citep{nureg1537p1} &
  Parte 2, \S1.2, \emph{Acceptance Criteria} (4 itens) \citep{nureg1537p2} &
  Tabela-resumo; tabela de ap\^endices \\
1.3 &
  Parte 1, \S1.3 \citep{nureg1537p1} &
  Parte 2, \S1.3, \emph{Acceptance Criteria} (4 itens) \citep{nureg1537p2} &
  Tabela-resumo; tabela de ap\^endices \\
1.4 &
  Parte 1, \S1.4 \citep{nureg1537p1} &
  Parte 2, \S1.4, \emph{Acceptance Criteria} (3 itens) \citep{nureg1537p2} &
  Tabela-resumo; tabela de ap\^endices \\
1.5 &
  Parte 1, \S1.5 \citep{nureg1537p1} &
  Parte 2, \S1.5, \emph{Acceptance Criteria} (2 itens) \citep{nureg1537p2} &
  Tabela-resumo; tabela de ap\^endices; substitui\c{c}\~ao de ``NRC/DOE''
  por ``[autoridade reguladora nacional] / organiza\c{c}\~oes de
  pesquisa reconhecidas'' \\
1.6 &
  Parte 1, \S1.6 \citep{nureg1537p1} &
  Parte 2, \S1.6, \emph{Acceptance Criteria} (3 itens) \citep{nureg1537p2} &
  Tabela-resumo; tabela de ap\^endices; refer\^encia expl\'icita ao
  Cap\'itulo~14 (Especifica\c{c}\~oes T\'ecnicas), impl\'icita no
  original \\
1.7 &
  Parte 1, \S1.7 \citep{nureg1537p1} (espec\'ifica dos EUA -- ver
  ressalva no pr\'oprio texto) &
  Parte 2, \S1.7, \emph{Acceptance Criteria} (2 itens) \citep{nureg1537p2};
  3\textordmasculine\ item ``fora dos EUA'' \'e adapta\c{c}\~ao pr\'opria &
  Tabela-resumo; tabela de ap\^endices; bloco ``fora dos Estados
  Unidos'' (sem equivalente no original) \\
1.8 &
  Parte 1, \S1.8 \citep{nureg1537p1} &
  Parte 2, \S1.8, \emph{Acceptance Criterion} (1 item, gen\'erico)
  \citep{nureg1537p2} &
  Tabela-resumo; tabela de ap\^endices; distin\c{c}\~ao expl\'icita entre
  instala\c{c}\~ao nova/renova\c{c}\~ao (impl\'icita no original) \\
\end{tabelametodologia}

\nusecbreak{anexo2}{Anexo}
\section*{Anexo -- Estrat\'egia de Licenciamento Faseado (Escopo Inicial
Limitado ao Pr\'edio do Reator)}
\addcontentsline{toc}{section}{Anexo -- Estrat\'egia de Licenciamento
Faseado (Escopo Inicial Limitado ao Pr\'edio do Reator)}
\label{sec:anexo-fase}

Este anexo registra a estrat\'egia discutida para limitar o escopo
inicial deste RPAS ao pr\'edio do reator e suas interfaces, deixando as
instala\c{c}\~oes de irradia\c{c}\~ao/experimentais associadas para uma
fase posterior de licenciamento.

\textbf{A NUREG-1537 n\~ao prev\^e um mecanismo formal de licenciamento
faseado.} O conceito de \emph{Interface Requirements} -- declara\c{c}\~ao
formal de requisitos para partes do projeto ainda n\~ao finalizadas -- \'e
espec\'ifico da Certifica\c{c}\~ao de Projeto sob 10~CFR Parte~52
\citep{cfr10p52}, via que este RPAS n\~ao segue (ver Se\c{c}\~oes~1.1 a
1.8, enquadradas em permiss\~ao de constru\c{c}\~ao/licen\c{c}a de
opera\c{c}\~ao sob 10~CFR Parte~50). Sob a NUREG-1537, o instrumento
dispon\'ivel para limitar o escopo inicial \'e a Se\c{c}\~ao~1.4,
``Instala\c{c}\~oes e Equipamentos Compartilhados'', combinada com o
mecanismo de emenda de licen\c{c}a da Se\c{c}\~ao~1.8.

\begin{itemize}
\item \textbf{Declarar o escopo.} Na Se\c{c}\~ao~1.4, declarar
  explicitamente que as instala\c{c}\~oes de irradia\c{c}\~ao n\~ao fazem
  parte deste RPAS/desta licen\c{c}a inicial -- a pr\'opria Parte~1
  \citep{nureg1537p1} prev\^e a declara\c{c}\~ao expl\'icita de aus\^encia
  de instala\c{c}\~oes compartilhadas; o caminho inverso, declarar o que
  fica de fora do escopo, segue a mesma l\'ogica.
\item \textbf{Descrever a interface, n\~ao a instala\c{c}\~ao.} Os dois
  crit\'erios de aceita\c{c}\~ao da Se\c{c}\~ao~1.4 (ver tabela de
  requisitos de aceita\c{c}\~ao daquela se\c{c}\~ao) exigem que o pr\'edio
  do reator seja projetado para acomodar qualquer uso ou falha das
  instala\c{c}\~oes compartilhadas sem degrada\c{c}\~ao das
  caracter\'isticas de seguran\c{c}a do reator, e que o projeto evite a
  propaga\c{c}\~ao de contamina\c{c}\~ao entre eles. Devem ser descritas
  as barreiras f\'isicas, el\'etricas, de HVAC e de
  confinamento/conten\c{c}\~ao entre o pr\'edio do reator e as futuras
  instala\c{c}\~oes de irradia\c{c}\~ao -- n\~ao o conte\'udo interno
  delas.
\item \textbf{Tratar provis\~oes f\'isicas j\'a embutidas como
  ``reservadas''.} Caso o n\'ucleo/vaso do reator j\'a nasça com canais de
  irradia\c{c}\~ao \emph{in-core} ou \emph{in-reflector} vazios (para uso
  futuro), isso deve constar estruturalmente no Cap\'itulo~4
  (descri\c{c}\~ao do reator), mas o Cap\'itulo~10, ``Experimental
  Facilities and Utilization'' \citep{nureg1537p1} -- que \'e onde a
  NUREG-1537 trata integralmente as instala\c{c}\~oes experimentais e de
  irradia\c{c}\~ao, incluindo incore facilities, in-reflector facilities,
  automatic transfer (``rabbit'') facilities e beam ports -- pode
  trat\'a-las como ``reservadas, sem uso experimental nesta fase'', sem
  exigir a an\'alise de seguran\c{c}a completa de experimentos que ainda
  n\~ao existem. Note-se que a pr\'opria NUREG-1537 registra que a falha
  de um experimento pode constituir o MHA do reator -- raz\~ao adicional
  para não antecipar an\'alises de experimentos ainda indefinidos.
\item \textbf{Usar a Se\c{c}\~ao~1.8/10~CFR~50.59 como porta de entrada
  para a expans\~ao futura.} Quando as instala\c{c}\~oes de
  irradia\c{c}\~ao forem de fato constru\'idas, sua incorpora\c{c}\~ao se
  d\'a por emenda de licen\c{c}a -- ou, quando a mudan\c{c}a n\~ao
  envolver quest\~ao de seguran\c{c}a n\~ao revisada, pelo processo
  autoimplementado equivalente \`a Se\c{c}\~ao~50.59 do 10~CFR~50
  \citep{cfr10p50}, j\'a referenciado na Se\c{c}\~ao~1.8 deste cap\'itulo.
\end{itemize}

\begin{notaregulatoria}{Estrat\'egia de redação, n\~ao conceito da NUREG-1537}
N\~ao h\'a, na NUREG-1537, um nome formal para essa estrat\'egia --
diferentemente dos \emph{Interface Requirements} de 10~CFR Parte~52. O
que se descreve aqui \'e o uso combinado das Se\c{c}\~oes~1.4 e~1.8 (e do
Cap\'itulo~10 para o detalhamento das instala\c{c}\~oes de
irradia\c{c}\~ao quando incorporadas) como fronteira de escopo declarada,
n\~ao um mecanismo regulat\'orio dedicado. Cabe ao requerente confirmar
com o \'org\~ao regulador nacional se essa abordagem \'e aceit\'avel para
o licenciamento faseado pretendido.
\end{notaregulatoria}
